% Options for packages loaded elsewhere
\PassOptionsToPackage{unicode}{hyperref}
\PassOptionsToPackage{hyphens}{url}
\documentclass[
]{ltjsarticle}
\usepackage{xcolor}
\usepackage[margin=2cm]{geometry}
\usepackage{amsmath,amssymb}
\setcounter{secnumdepth}{-\maxdimen} % remove section numbering
\usepackage{iftex}
\ifPDFTeX
  \usepackage[T1]{fontenc}
  \usepackage[utf8]{inputenc}
  \usepackage{textcomp} % provide euro and other symbols
\else % if luatex or xetex
  \usepackage{unicode-math} % this also loads fontspec
  \defaultfontfeatures{Scale=MatchLowercase}
  \defaultfontfeatures[\rmfamily]{Ligatures=TeX,Scale=1}
\fi
\usepackage{lmodern}
\ifPDFTeX\else
  % xetex/luatex font selection
  \setmainfont[]{Hiragino Sans}
  \setmonofont[]{Menlo}
\fi
% Use upquote if available, for straight quotes in verbatim environments
\IfFileExists{upquote.sty}{\usepackage{upquote}}{}
\IfFileExists{microtype.sty}{% use microtype if available
  \usepackage[]{microtype}
  \UseMicrotypeSet[protrusion]{basicmath} % disable protrusion for tt fonts
}{}
\makeatletter
\@ifundefined{KOMAClassName}{% if non-KOMA class
  \IfFileExists{parskip.sty}{%
    \usepackage{parskip}
  }{% else
    \setlength{\parindent}{0pt}
    \setlength{\parskip}{6pt plus 2pt minus 1pt}}
}{% if KOMA class
  \KOMAoptions{parskip=half}}
\makeatother
\usepackage{color}
\usepackage{fancyvrb}
\newcommand{\VerbBar}{|}
\newcommand{\VERB}{\Verb[commandchars=\\\{\}]}
\DefineVerbatimEnvironment{Highlighting}{Verbatim}{commandchars=\\\{\}}
% Add ',fontsize=\small' for more characters per line
\newenvironment{Shaded}{}{}
\newcommand{\AlertTok}[1]{\textcolor[rgb]{1.00,0.00,0.00}{\textbf{#1}}}
\newcommand{\AnnotationTok}[1]{\textcolor[rgb]{0.38,0.63,0.69}{\textbf{\textit{#1}}}}
\newcommand{\AttributeTok}[1]{\textcolor[rgb]{0.49,0.56,0.16}{#1}}
\newcommand{\BaseNTok}[1]{\textcolor[rgb]{0.25,0.63,0.44}{#1}}
\newcommand{\BuiltInTok}[1]{\textcolor[rgb]{0.00,0.50,0.00}{#1}}
\newcommand{\CharTok}[1]{\textcolor[rgb]{0.25,0.44,0.63}{#1}}
\newcommand{\CommentTok}[1]{\textcolor[rgb]{0.38,0.63,0.69}{\textit{#1}}}
\newcommand{\CommentVarTok}[1]{\textcolor[rgb]{0.38,0.63,0.69}{\textbf{\textit{#1}}}}
\newcommand{\ConstantTok}[1]{\textcolor[rgb]{0.53,0.00,0.00}{#1}}
\newcommand{\ControlFlowTok}[1]{\textcolor[rgb]{0.00,0.44,0.13}{\textbf{#1}}}
\newcommand{\DataTypeTok}[1]{\textcolor[rgb]{0.56,0.13,0.00}{#1}}
\newcommand{\DecValTok}[1]{\textcolor[rgb]{0.25,0.63,0.44}{#1}}
\newcommand{\DocumentationTok}[1]{\textcolor[rgb]{0.73,0.13,0.13}{\textit{#1}}}
\newcommand{\ErrorTok}[1]{\textcolor[rgb]{1.00,0.00,0.00}{\textbf{#1}}}
\newcommand{\ExtensionTok}[1]{#1}
\newcommand{\FloatTok}[1]{\textcolor[rgb]{0.25,0.63,0.44}{#1}}
\newcommand{\FunctionTok}[1]{\textcolor[rgb]{0.02,0.16,0.49}{#1}}
\newcommand{\ImportTok}[1]{\textcolor[rgb]{0.00,0.50,0.00}{\textbf{#1}}}
\newcommand{\InformationTok}[1]{\textcolor[rgb]{0.38,0.63,0.69}{\textbf{\textit{#1}}}}
\newcommand{\KeywordTok}[1]{\textcolor[rgb]{0.00,0.44,0.13}{\textbf{#1}}}
\newcommand{\NormalTok}[1]{#1}
\newcommand{\OperatorTok}[1]{\textcolor[rgb]{0.40,0.40,0.40}{#1}}
\newcommand{\OtherTok}[1]{\textcolor[rgb]{0.00,0.44,0.13}{#1}}
\newcommand{\PreprocessorTok}[1]{\textcolor[rgb]{0.74,0.48,0.00}{#1}}
\newcommand{\RegionMarkerTok}[1]{#1}
\newcommand{\SpecialCharTok}[1]{\textcolor[rgb]{0.25,0.44,0.63}{#1}}
\newcommand{\SpecialStringTok}[1]{\textcolor[rgb]{0.73,0.40,0.53}{#1}}
\newcommand{\StringTok}[1]{\textcolor[rgb]{0.25,0.44,0.63}{#1}}
\newcommand{\VariableTok}[1]{\textcolor[rgb]{0.10,0.09,0.49}{#1}}
\newcommand{\VerbatimStringTok}[1]{\textcolor[rgb]{0.25,0.44,0.63}{#1}}
\newcommand{\WarningTok}[1]{\textcolor[rgb]{0.38,0.63,0.69}{\textbf{\textit{#1}}}}
\usepackage{longtable,booktabs,array}
\newcounter{none} % for unnumbered tables
\usepackage{calc} % for calculating minipage widths
% Correct order of tables after \paragraph or \subparagraph
\usepackage{etoolbox}
\makeatletter
\patchcmd\longtable{\par}{\if@noskipsec\mbox{}\fi\par}{}{}
\makeatother
% Allow footnotes in longtable head/foot
\IfFileExists{footnotehyper.sty}{\usepackage{footnotehyper}}{\usepackage{footnote}}
\makesavenoteenv{longtable}
\setlength{\emergencystretch}{3em} % prevent overfull lines
\providecommand{\tightlist}{%
  \setlength{\itemsep}{0pt}\setlength{\parskip}{0pt}}
\usepackage{bookmark}
\IfFileExists{xurl.sty}{\usepackage{xurl}}{} % add URL line breaks if available
\urlstyle{same}
\hypersetup{
  hidelinks,
  pdfcreator={LaTeX via pandoc}}

\author{}
\date{}

\begin{document}

{
\setcounter{tocdepth}{2}
\tableofcontents
}
\section{進化計算による「LLM 駆動 MAP-Elites を安全に走らせる言語」探索
---
実験報告}\label{ux9032ux5316ux8a08ux7b97ux306bux3088ux308bllm-ux99c6ux52d5-map-elites-ux3092ux5b89ux5168ux306bux8d70ux3089ux305bux308bux8a00ux8a9eux63a2ux7d22-ux5b9fux9a13ux5831ux544a}

\textbf{日付:} 2026-05-16 \textbf{実験名:}
\texttt{2026-05-16\_robustness\_v1} \textbf{作業ディレクトリ:}
\texttt{/Users/kodamay/ocaml-app/abclcp-project/aice-pi-evolution/experiments/2026-05-16\_robustness\_v1/}
\textbf{親実験:} \href{../../REPORT_JA.md}{REPORT\_JA.md} (PiLang
を生んだ 2026-05-15 実験)

\begin{center}\rule{0.5\linewidth}{0.5pt}\end{center}

\subsection{1. 目的}\label{ux76eeux7684}

PiLang を生んだ進化計算は \textbf{「π 10,000 桁を求めるのに最適な言語」}
を探索した。本実験は同じパイプラインを\textbf{メタ層}に当て、

\begin{quote}
\textbf{「LLM 駆動 MAP-Elites を実機上で安全 (コスト ≤ \$2 / 壁時計 ≤ 30
分 / silent death 0) に完走させる」}
\end{quote}

という問題を解くのに最適な\textbf{ランタイム構成言語}を進化計算で探索する。元の
PiLang
探索が「数値計算問題から言語を生む」だったのに対し、本実験は「メタ計算問題から言語を生む」点で対照的である。

具体的には、2026-05-16 セッションの前半で観測した 3 つの失敗モード (試行
\#1〜\#3) を踏まえ、\texttt{.aice\ (YAML)\ →\ .ga.json\ →\ .aipl}
のパイプラインを通して、

\begin{enumerate}
\def\labelenumi{\arabic{enumi}.}
\tightlist
\item
  11 軸 (provider, model, concurrency, seed\_count, generations,
  request\_timeout\_s, sdk\_retries, drain\_timeout\_s, idle\_ms,
  checkpoint\_per\_gen, caffeinate) を持つ
  \textbf{ランタイム遺伝子}を設計、
\item
  F1〜F5 の 5 つの reviewer を LLM-as-judge として実装、
\item
  OpenAI gpt-4o-mini で MAP-Elites を 30 世代走らせ、
\item
  個体群から「安全に走るための設計原理」を抽出、
\item
  それを表現する言語 \textbf{PhiLang} を提案する、
\end{enumerate}

までを一気通貫で行った。

\subsection{2. 実験構成}\label{ux5b9fux9a13ux69cbux6210}

\subsubsection{2.1
ディレクトリ}\label{ux30c7ux30a3ux30ecux30afux30c8ux30ea}

\begin{verbatim}
experiments/2026-05-16_robustness_v1/
├── Pi_Phase1_Robustness.aice          ← 案 1 (YAML) の問題定義
├── Pi_Phase1_Robustness.math.aice     ← 案 3 (数理最適化) の問題定義 (参考)
├── Pi_Phase1_Robustness.ga.json       ← 中間 IR
├── Pi_Phase1_Robustness.aipl          ← 実行プログラム (21 KB)
├── robustness_paradigm.schema.json    ← 遺伝子スキーマ
├── Pi_Phase1_Robustness.aipl_lineage.json   ← 38 個体の結果
├── ai_usage.json                      ← OpenAI 課金
├── Pi_Phase1_Robustness.log           ← 実行ログ
└── REPORT.md                          ← この文書
\end{verbatim}

\subsubsection{\texorpdfstring{2.2 新規遺伝子スキーマ
\texttt{robustness\_paradigm.schema.json}}{2.2 新規遺伝子スキーマ robustness\_paradigm.schema.json}}\label{ux65b0ux898fux907aux4f1dux5b50ux30b9ux30adux30fcux30de-robustness_paradigm.schema.json}

PiLang スキーマ (\texttt{pi\_paradigm.schema.json})
を\textbf{継承せず}、メタ層用に独立に設計した。11 軸:

{\def\LTcaptype{none} % do not increment counter
\begin{longtable}[]{@{}
  >{\raggedright\arraybackslash}p{(\linewidth - 4\tabcolsep) * \real{0.3333}}
  >{\raggedright\arraybackslash}p{(\linewidth - 4\tabcolsep) * \real{0.3333}}
  >{\raggedright\arraybackslash}p{(\linewidth - 4\tabcolsep) * \real{0.3333}}@{}}
\toprule\noalign{}
\begin{minipage}[b]{\linewidth}\raggedright
軸
\end{minipage} & \begin{minipage}[b]{\linewidth}\raggedright
値域
\end{minipage} & \begin{minipage}[b]{\linewidth}\raggedright
由来
\end{minipage} \\
\midrule\noalign{}
\endhead
\bottomrule\noalign{}
\endlastfoot
\texttt{provider} & \texttt{openai,\ anthropic,\ vertex,\ mock} & LLM
接続先 \\
\texttt{model} & \texttt{gpt-4o-mini,\ gpt-4o,\ claude-haiku-4-5} &
モデル \\
\texttt{concurrency} & \texttt{c1,\ c2,\ c4,\ c8,\ c16,\ c32} &
\texttt{ABCL\_AI\_MAX\_CONCURRENT} \\
\texttt{seed\_count} & \texttt{s4,\ s6,\ s8,\ s12,\ s16,\ s24} &
MAP-Elites 種数 \\
\texttt{generations} & \texttt{g10,\ g15,\ g20,\ g30,\ g40,\ g60,\ g80}
& MAP-Elites 世代数 \\
\texttt{request\_timeout\_s} & \texttt{t10,\ t30,\ t60,\ t120,\ t300} &
F2 (\texttt{AIPL\_AI\_REQUEST\_TIMEOUT}) \\
\texttt{sdk\_retries} & \texttt{r0,\ r1,\ r3,\ r5} & F2
(\texttt{AIPL\_AI\_SDK\_RETRIES}) \\
\texttt{drain\_timeout\_s} &
\texttt{d60,\ d600,\ d3600,\ d7200,\ d14400} & F1
(\texttt{-\/-timeout}) \\
\texttt{idle\_ms} & \texttt{i1000,\ i5000,\ i30000,\ i60000,\ i300000} &
F1 (\texttt{-\/-idle-ms}) \\
\texttt{checkpoint\_per\_gen} & \texttt{off,\ on} & F3 (lineage dump
頻度) \\
\texttt{caffeinate} & \texttt{off,\ on} & F5 (Mac sleep 対策) \\
\end{longtable}
}

7 軸を \texttt{ordinal\_axes} で順序付け、6 つの \texttt{coherence}
ルール (例:
\texttt{provider=openai\ ⇒\ drain\_timeout\_s\ ∈\ \{d3600,\ d7200,\ d14400\}})
で「明らかに不安全な組合せ」を変異が生まないようにした。\texttt{open\_axes=false}
--- ランタイム設定なので LLM が新軸を提案する余地はない。

\subsubsection{2.3 失敗モード F1〜F5 を 5 reviewer
に展開}\label{ux5931ux6557ux30e2ux30fcux30c9-f1f5-ux3092-5-reviewer-ux306bux5c55ux958b}

2026-05-16 セッション前半で観測した失敗モードを 1 mode = 1 reviewer
に対応させた:

{\def\LTcaptype{none} % do not increment counter
\begin{longtable}[]{@{}
  >{\raggedright\arraybackslash}p{(\linewidth - 6\tabcolsep) * \real{0.2308}}
  >{\raggedright\arraybackslash}p{(\linewidth - 6\tabcolsep) * \real{0.2308}}
  >{\raggedleft\arraybackslash}p{(\linewidth - 6\tabcolsep) * \real{0.3077}}
  >{\raggedright\arraybackslash}p{(\linewidth - 6\tabcolsep) * \real{0.2308}}@{}}
\toprule\noalign{}
\begin{minipage}[b]{\linewidth}\raggedright
reviewer
\end{minipage} & \begin{minipage}[b]{\linewidth}\raggedright
検出対象
\end{minipage} & \begin{minipage}[b]{\linewidth}\raggedleft
weight
\end{minipage} & \begin{minipage}[b]{\linewidth}\raggedright
スコア基準
\end{minipage} \\
\midrule\noalign{}
\endhead
\bottomrule\noalign{}
\endlastfoot
F1\_IdleTimeoutGuard & \texttt{drain\_timeout\_s} と \texttt{idle\_ms}
が real LLM 用に十分か & 0.20 & \texttt{d3600+\ \&\ i30000+} → 1.0, 片方
→ 0.5, 両方未満 → 0.0 \\
F2\_SdkHangGuard & \texttt{request\_timeout\_s} と \texttt{sdk\_retries}
で SDK ハング回避 & 0.25 & \texttt{t30..t120\ \&\ r3+} → 1.0,
\texttt{t10..t300\ \&\ r1+} → 0.5 \\
F3\_CheckpointGuard & \texttt{checkpoint\_per\_gen=on} か & 0.15 &
\texttt{on\ →\ 1.0,\ off\ →\ 0.0} (二値) \\
F4\_ThroughputGuard & 30 分以内に完走可能か
(\texttt{total\_calls\ /\ (23·max(1,α·c))}) & 0.25 &
\texttt{wall\ ≤\ 15min} → 1.0, \ldots, \texttt{\textgreater{}\ 60min} →
0.0 \\
F5\_SleepGuard & wall \textgreater{} 30min なら \texttt{caffeinate=on}
か & 0.15 & 推定時間 ≤ 30min → 1.0, \textgreater{} 30min かつ on → 1.0,
それ以外 → 0.0 \\
\end{longtable}
}

各 persona には \textbf{prefix エンコーディングキー}
(\texttt{drain\_timeout\_s\ の\ d\{N\}\ は\ N\ 秒,\ d3600=1時間}) と
\textbf{0.0--1.0 段階スコア基準} を明記。これは smoke 縮小版 (4 seed ×
10 gen) で reviewer が OK/NG 二値を返そうとして全スコアが 0.04〜0.23
に張り付いた問題を直したもの。

\subsubsection{\texorpdfstring{2.4 \texttt{.aice} から \texttt{.aipl}
までのパイプライン}{2.4 .aice から .aipl までのパイプライン}}\label{aice-ux304bux3089-.aipl-ux307eux3067ux306eux30d1ux30a4ux30d7ux30e9ux30a4ux30f3}

PiLang 実験で確立した \texttt{.aice\ →\ .ga.json\ →\ .aipl}
の三段パイプラインをそのまま流用。今回は \texttt{.aice} の dialect が
\textbf{YAML (案 1)} であるため、\texttt{aice\_parser.lower()} の
curly-brace パーサは使えず、\textbf{YAML を手作業で \texttt{.ga.json}
にマップ}した。\texttt{.ga.json\ →\ .aipl} は既存の
\texttt{aipl\_codegen.generate\_program()} がそのまま動いた (641 行の
AIPL プログラムを生成)。

\subsection{3.
パイプライン実行}\label{ux30d1ux30a4ux30d7ux30e9ux30a4ux30f3ux5b9fux884c}

\subsubsection{3.1 試行履歴 (3 失敗 + 1
完走)}\label{ux8a66ux884cux5c65ux6b74-3-ux5931ux6557-1-ux5b8cux8d70}

{\def\LTcaptype{none} % do not increment counter
\begin{longtable}[]{@{}
  >{\raggedright\arraybackslash}p{(\linewidth - 4\tabcolsep) * \real{0.3333}}
  >{\raggedright\arraybackslash}p{(\linewidth - 4\tabcolsep) * \real{0.3333}}
  >{\raggedright\arraybackslash}p{(\linewidth - 4\tabcolsep) * \real{0.3333}}@{}}
\toprule\noalign{}
\begin{minipage}[b]{\linewidth}\raggedright
試行
\end{minipage} & \begin{minipage}[b]{\linewidth}\raggedright
結果
\end{minipage} & \begin{minipage}[b]{\linewidth}\raggedright
教訓
\end{minipage} \\
\midrule\noalign{}
\endhead
\bottomrule\noalign{}
\endlastfoot
\#1 & 5 秒で死亡 (1 call) & \texttt{-\/-timeout\ 2.0} /
\texttt{-\/-idle-ms\ 120} の既定値が real LLM では短すぎ \\
\#2 & gen 38/40 (1267 calls, \$0.44) で silent death & OpenAI SDK の
HTTPS request にタイムアウト無で永久待機 \\
\#3 & gen 11/40 (23 個体, \$0.20) でユーザ中断 & 実効並列度 1〜2
でスループット遅い \\
smoke & 14 個体 (350 calls, \$0.094, 13 分) 完走 & reviewer OK/NG 表記で
score 低 (max 0.23) \\
\textbf{\#4 本番} & \textbf{38 個体 (950 calls, \$0.269, 33 分) 完走} &
reviewer 0.0--1.0 + prefix キー化で score 改善 (max 0.46) \\
\end{longtable}
}

\subsubsection{3.2 \#4
本番ランの安全策}\label{ux672cux756aux30e9ux30f3ux306eux5b89ux5168ux7b56}

\begin{Shaded}
\begin{Highlighting}[]
\VariableTok{AIPL\_AI\_PROVIDER}\OperatorTok{=}\NormalTok{openai}
\VariableTok{ABCL\_AI\_MAX\_CONCURRENT}\OperatorTok{=}\NormalTok{8}
\VariableTok{AIPL\_AI\_REQUEST\_TIMEOUT}\OperatorTok{=}\NormalTok{60      }\CommentTok{\# F2 対策}
\VariableTok{AIPL\_AI\_SDK\_RETRIES}\OperatorTok{=}\NormalTok{3           }\CommentTok{\# F2 対策}
\ExtensionTok{{-}{-}timeout}\NormalTok{ 7200                  }\CommentTok{\# F1 対策}
\ExtensionTok{{-}{-}idle{-}ms}\NormalTok{ 60000                 }\CommentTok{\# F1 対策}
\end{Highlighting}
\end{Shaded}

これら 4 つの env と CLI フラグは試行 \#1〜\#3
から学んだ「失敗を防ぐ既定」である。\textbf{興味深いことに、進化計算は同じ結論にメタ的に到達した}
(§4 参照)。

\subsubsection{3.3 サマリ}\label{ux30b5ux30deux30ea}

{\def\LTcaptype{none} % do not increment counter
\begin{longtable}[]{@{}lrrl@{}}
\toprule\noalign{}
指標 & 実測 & 予測 & 差 \\
\midrule\noalign{}
\endhead
\bottomrule\noalign{}
\endlastfoot
完走 & ✓ & --- & --- \\
壁時計 & 33 分 & 45 分 & −27\% \\
コスト & \$0.269 & \$0.34 & −20\% \\
コール数 & 950 & 1,026 & −7\% \\
スループット & 30 calls/min & 23 calls/min & +30\% \\
silent\_death & 0 & 0 & ✓ \\
\end{longtable}
}

\subsection{4. 結果}\label{ux7d50ux679c}

\subsubsection{4.1 個体軸頻度 (38 個体, top
3)}\label{ux500bux4f53ux8ef8ux983bux5ea6-38-ux500bux4f53-top-3}

{\def\LTcaptype{none} % do not increment counter
\begin{longtable}[]{@{}
  >{\raggedright\arraybackslash}p{(\linewidth - 4\tabcolsep) * \real{0.3333}}
  >{\raggedright\arraybackslash}p{(\linewidth - 4\tabcolsep) * \real{0.3333}}
  >{\raggedright\arraybackslash}p{(\linewidth - 4\tabcolsep) * \real{0.3333}}@{}}
\toprule\noalign{}
\begin{minipage}[b]{\linewidth}\raggedright
軸
\end{minipage} & \begin{minipage}[b]{\linewidth}\raggedright
上位値
\end{minipage} & \begin{minipage}[b]{\linewidth}\raggedright
コメント
\end{minipage} \\
\midrule\noalign{}
\endhead
\bottomrule\noalign{}
\endlastfoot
\texttt{provider} & openai:14, mock:13, vertex:11 & 偏り少
(rng\_seed=31415926 起因の偶然均一) \\
\texttt{model} & gpt-4o:15, claude-haiku-4-5:14, gpt-4o-mini:9 &
reviewer は model ID をテキスト判定するだけなので gpt-4o-mini
が必ずしも勝たない \\
\texttt{concurrency} & c8:9, c16:9, c1:8, c32:7, c2:5 & top 5 では c1
が支持 (g15 だと並列度に依存しない) \\
\texttt{seed\_count} & s8:13, s6:8, s12:8, s4:5, s16:4 &
中央値が好まれた \\
\texttt{generations} & \textbf{g15:14}, g10:6, g20:5, g30:4, g40:4 &
\textbf{「短いほど完走確率高い」が圧倒的多数派} \\
\texttt{request\_timeout\_s} & t60:11, t30:10, t120:8, t300:5, t10:4 &
F2 推奨域 \texttt{t30-t120} に集中 \\
\texttt{sdk\_retries} & r3:12, r5:11, r1:8, r0:7 & \textbf{r3+ が 23/38
(60\%)} で F2 推奨を満たす \\
\texttt{drain\_timeout\_s} & d7200:11, d14400:8, d3600:7, d600:7, d60:5
& \textbf{d3600+ が 26/38 (68\%)} で F1 推奨を満たす \\
\texttt{idle\_ms} & i60000:11, i30000:10, i300000:8, i5000:5, i1000:4 &
\textbf{i30000+ が 29/38 (76\%)} で F1 推奨を満たす \\
\texttt{checkpoint\_per\_gen} & off:22, on:16 & F3 weight=0.15
が他に押し負け \\
\texttt{caffeinate} & off:20, on:18 & ほぼ五分 \\
\end{longtable}
}

→ \textbf{F1 と F2 の reviewer が強い選択圧を生み、F3 は weight
不足で機能不全}、という非対称な学習結果。

\subsubsection{4.2 Top 5
個体の遺伝子}\label{top-5-ux500bux4f53ux306eux907aux4f1dux5b50}

{\def\LTcaptype{none} % do not increment counter
\begin{longtable}[]{@{}
  >{\raggedleft\arraybackslash}p{(\linewidth - 10\tabcolsep) * \real{0.1905}}
  >{\raggedright\arraybackslash}p{(\linewidth - 10\tabcolsep) * \real{0.1429}}
  >{\raggedleft\arraybackslash}p{(\linewidth - 10\tabcolsep) * \real{0.1905}}
  >{\raggedright\arraybackslash}p{(\linewidth - 10\tabcolsep) * \real{0.1429}}
  >{\raggedleft\arraybackslash}p{(\linewidth - 10\tabcolsep) * \real{0.1905}}
  >{\raggedright\arraybackslash}p{(\linewidth - 10\tabcolsep) * \real{0.1429}}@{}}
\toprule\noalign{}
\begin{minipage}[b]{\linewidth}\raggedleft
順位
\end{minipage} & \begin{minipage}[b]{\linewidth}\raggedright
id
\end{minipage} & \begin{minipage}[b]{\linewidth}\raggedleft
gen
\end{minipage} & \begin{minipage}[b]{\linewidth}\raggedright
op
\end{minipage} & \begin{minipage}[b]{\linewidth}\raggedleft
score
\end{minipage} & \begin{minipage}[b]{\linewidth}\raggedright
genome (要約)
\end{minipage} \\
\midrule\noalign{}
\endhead
\bottomrule\noalign{}
\endlastfoot
1 & I25 & 17 & axis\_resample & \textbf{0.457} & mock·gpt-4o · c1·s6·g15
· t30·r3·d14400·i30000 · ckpt=off·caff=off \\
2 & I24 & 16 & uniform\_crossover & 0.455 & mock·gpt-4o · c8·s8·g15 ·
t300·r5·d600·i30000 · ckpt=off·caff=on \\
3 & I19 & 11 & uniform\_crossover & 0.421 & openai·claude-haiku ·
c2·s8·g15 · t30·r5·d7200·i60000 · ckpt=off·caff=off \\
4 & I1 & 0 & seed & 0.397 & mock·gpt-4o · c1·s6·g15 · t30·r3·d600·i30000
· ckpt=off·caff=off \\
5 & I30 & 22 & axis\_resample & 0.396 & openai·claude-haiku · c32·s8·g15
· t30·r5·d7200·i60000 · ckpt=off·caff=off \\
\end{longtable}
}

\subsubsection{4.3
世代別スコア推移}\label{ux4e16ux4ee3ux5225ux30b9ux30b3ux30a2ux63a8ux79fb}

\begin{verbatim}
gen= 0: max=0.397 mean=0.257   (8 seed の初期スコア帯)
gen= 4: max=0.392
gen=11: max=0.421              (crossover 1 発目で上振れ)
gen=16: max=0.455
gen=17: max=0.457              ← 全体ベスト
gen=22: max=0.396
\end{verbatim}

→ 25 世代で 0.397 → 0.457 に \textbf{約 15\%
改善}。改善幅は小さいがブレイクスルー (gen 16-17) は crossover
が生み出した。MAP-Elites の elite
更新と突然変異の組合せが機能していることが分かる。

\subsection{5.
進化計算が示唆する設計原理}\label{ux9032ux5316ux8a08ux7b97ux304cux793aux5506ux3059ux308bux8a2dux8a08ux539fux7406}

38 個体のスコア分布と top 5 の共通項から、「LLM 駆動 MAP-Elites
を安全に走らせる」言語が\textbf{先天的に備えるべき四つの特性}を抽出する。これは
PiLang における
\texttt{Array\_DSL\ ×\ refinement-typed\ ×\ region-owned\ ×\ streaming}
の 4 つ巴と同じ抽象度の知見である。

{\def\LTcaptype{none} % do not increment counter
\begin{longtable}[]{@{}
  >{\raggedright\arraybackslash}p{(\linewidth - 2\tabcolsep) * \real{0.5000}}
  >{\raggedright\arraybackslash}p{(\linewidth - 2\tabcolsep) * \real{0.5000}}@{}}
\toprule\noalign{}
\begin{minipage}[b]{\linewidth}\raggedright
特性
\end{minipage} & \begin{minipage}[b]{\linewidth}\raggedright
進化計算が示した根拠
\end{minipage} \\
\midrule\noalign{}
\endhead
\bottomrule\noalign{}
\endlastfoot
\textbf{(A) Bounded burst} & top 5 全員が \texttt{generations=g15}
を支持。長時間ランより短バースト+部分結果採用を選好 \\
\textbf{(B) Mandatory retry/timeout} & 全 top 5 が
\texttt{sdk\_retries\ ≥\ r3} かつ
\texttt{request\_timeout\_s\ ≤\ t300}。「待ち続ける」より「諦めて再試行」 \\
\textbf{(C) Generous drain buffer} & 80\% が
\texttt{drain\_timeout\_s\ ≥\ d3600} \&
\texttt{idle\_ms\ ≥\ i30000}。actor 落穂拾いに余裕を持つ \\
\textbf{(D) Forced checkpoint} & 進化計算自体は F3
を見落としたが、人間が安全運用のために\textbf{必須化}すべき特性 (§7.4
参照) \\
\end{longtable}
}

→ \textbf{「Bounded burst × Mandatory retry × Generous drain × Forced
checkpoint」} の交点として、メタ問題に最適な言語が浮かび上がる。

\subsection{6. 進化計算が生んだ言語
PhiLang}\label{ux9032ux5316ux8a08ux7b97ux304cux751fux3093ux3060ux8a00ux8a9e-philang}

PiLang が \textbf{πの 1 タスク}から生まれた言語であるように、本実験は
\textbf{メタ・実行制御の 1 タスク}から \textbf{PhiLang} (Φ = Phase,
「位相」) を生んだ。

\subsubsection{6.1 PhiLang
のコア構文}\label{philang-ux306eux30b3ux30a2ux69cbux6587}

\begin{Shaded}
\begin{Highlighting}[]
\NormalTok{\# PhiLang — Phase Orchestration Language synthesized from MAP{-}Elites}
\NormalTok{\# (2026{-}05{-}16\_robustness\_v1 evolution result)}

\NormalTok{phase Phase1}
\NormalTok{  within   30 min      and   $0.30                     \# ハード予算}
\NormalTok{  uses     openai/gpt{-}4o{-}mini}
\NormalTok{  drains   for 2 h     idle 60 s                       \# actor drain buffer}
\NormalTok{  caffeine if duration \textgreater{} 30 min                        \# F5}
\NormalTok{\{}
\NormalTok{  seeds       = 8}
\NormalTok{  generations = 15                                     \# ← 短バースト (A)}
\NormalTok{  cell\_axes   = [paradigm, arithmetic, algorithm, precision, memory]}

\NormalTok{  every call has policy \{                              \# ← 全コールに強制}
\NormalTok{    timeout    : 60 s}
\NormalTok{    retries    : 3 with exponential\_backoff           \# ← (B)}
\NormalTok{    on\_hang    : kill\_and\_retry}
\NormalTok{  \}}

\NormalTok{  after every individual \{                             \# ← 強制 checkpoint (D)}
\NormalTok{    checkpoint lineage to "out/lineage.json"}
\NormalTok{  \}}
\NormalTok{\}}
\NormalTok{guarantees \{}
\NormalTok{  no\_silent\_death                                      \# ← (B), (C) から導出}
\NormalTok{  bounded\_cost            $0.30}
\NormalTok{  bounded\_wall            30 min}
\NormalTok{  partial\_resume\_on\_kill                               \# ← (D) から導出}
\NormalTok{\}}
\end{Highlighting}
\end{Shaded}

3 ブロック
(\texttt{phase\ /\ every\ call\ has\ policy\ /\ after\ every\ individual})
+ 1 契約ブロック (\texttt{guarantees})
で書ける。\texttt{every\ call\ has\ policy} は
\textbf{言語の意味論レベル}でレキシカルスコープ内の全 \texttt{ai\_call}
に強制適用される (オプトインではない)。

\subsubsection{6.2 脱糖形 (AIPL
ランタイム呼び出し)}\label{ux8131ux7cd6ux5f62-aipl-ux30e9ux30f3ux30bfux30a4ux30e0ux547cux3073ux51faux3057}

PhiLang の上記 1 ブロックは現状の AIPL
ランタイムに対しては次のように脱糖される:

\begin{verbatim}
# 環境変数 (every call has policy → SDK 設定)
ABCL_AI_MAX_CONCURRENT     = 8        # phase 内既定
AIPL_AI_REQUEST_TIMEOUT    = 60       # policy.timeout
AIPL_AI_SDK_RETRIES        = 3        # policy.retries
AIPL_AI_TOKEN_BUDGET       = 3_000_000  # within $0.30 から逆算

# CLI フラグ (drains 句)
--timeout  7200          # drains for 2 h
--idle-ms  60000         # idle 60 s

# AIPL コード (after every individual)
while (gen <= 15) do {
  ...                    # 個体生成と評価
  send lineage.add(...)
  send elite.propose(...)
  var n = now lineage.dump("out/lineage.json");   # ← 強制 checkpoint
  gen = gen + 1;
}

# シェルラッパ (caffeine 句)
caffeinate -i &          # PID をフェーズ終了で kill
\end{verbatim}

\subsubsection{6.3 4
つの言語要素が解決する失敗モード}\label{ux3064ux306eux8a00ux8a9eux8981ux7d20ux304cux89e3ux6c7aux3059ux308bux5931ux6557ux30e2ux30fcux30c9}

{\def\LTcaptype{none} % do not increment counter
\begin{longtable}[]{@{}
  >{\raggedright\arraybackslash}p{(\linewidth - 4\tabcolsep) * \real{0.3333}}
  >{\raggedright\arraybackslash}p{(\linewidth - 4\tabcolsep) * \real{0.3333}}
  >{\raggedright\arraybackslash}p{(\linewidth - 4\tabcolsep) * \real{0.3333}}@{}}
\toprule\noalign{}
\begin{minipage}[b]{\linewidth}\raggedright
PhiLang の構文要素
\end{minipage} & \begin{minipage}[b]{\linewidth}\raggedright
解決する失敗モード
\end{minipage} & \begin{minipage}[b]{\linewidth}\raggedright
進化計算での根拠
\end{minipage} \\
\midrule\noalign{}
\endhead
\bottomrule\noalign{}
\endlastfoot
\texttt{within\ 30\ min\ and\ \$0.30} (ハード予算) & F4
(low\_throughput) & top 5 全員が g15 → 短バースト原則 \\
\texttt{every\ call\ has\ policy\ \{\ timeout/retries\ \}} & F2
(sdk\_hang) & top 5 全員が \texttt{r3+} で SDK 強制リトライ \\
\texttt{drains\ for\ 2h\ idle\ 60s} & F1 (idle\_timeout\_too\_short) &
80\% が \texttt{d3600+}, \texttt{i30000+} \\
\texttt{after\ every\ individual\ \{\ checkpoint\ \}} & F3
(no\_checkpoint) & 進化計算が見落としたので\textbf{言語が必須化} \\
\texttt{caffeine\ if\ duration\ \textgreater{}\ 30min} & F5
(macos\_sleep) & reviewer F5 が学習済の判定式 \\
\texttt{guarantees\ \{\ ...\ \}} 節 & 全 F1〜F5
を\textbf{契約として明文化} & 個体群から導出した invariants \\
\end{longtable}
}

\subsubsection{6.4 PhiLang
の意味論的特徴}\label{philang-ux306eux610fux5473ux8ad6ux7684ux7279ux5fb4}

\paragraph{6.4.1
「コールサイト無し」のリトライ/タイムアウト}\label{ux30b3ux30fcux30ebux30b5ux30a4ux30c8ux7121ux3057ux306eux30eaux30c8ux30e9ux30a4ux30bfux30a4ux30e0ux30a2ux30a6ux30c8}

PhiLang では \texttt{ai\_call(...)}
一文一文にリトライやタイムアウトを書かない。\texttt{every\ call\ has\ policy}
ブロックが\textbf{レキシカル}に効くため、生 LLM API
のサブセットは存在しない。これは試行 \#2 で SDK
が暗黙にハングした問題を「言語の表現可能性」レベルで除去する。

\paragraph{6.4.2 「チェックポイント無し phase
は型エラー」}\label{ux30c1ux30a7ux30c3ux30afux30ddux30a4ux30f3ux30c8ux7121ux3057-phase-ux306fux578bux30a8ux30e9ux30fc}

PhiLang のフェーズは
\texttt{after\ every\ individual\ \{\ checkpoint\ \}}
句が必須。これを書き忘れると型検査で
\texttt{non-resumable\ phase:\ missing\ checkpoint\ clause} を出す。試行
\#2 で 38/40 個体分のスコアが silent death
で消えた問題を\textbf{コンパイル時}に防ぐ。

\paragraph{\texorpdfstring{6.4.3 \texttt{guarantees} 節は実行可能
invariant}{6.4.3 guarantees 節は実行可能 invariant}}\label{guarantees-ux7bc0ux306fux5b9fux884cux53efux80fd-invariant}

\texttt{bounded\_cost\ \$0.30} は単なるコメントではなく、ランタイムが
\texttt{ai\_usage.json} を監視して超過時に \texttt{BudgetExceeded}
を投げる。\texttt{no\_silent\_death} は
\texttt{every\ call\ has\ policy}
の存在から自動証明される。\textbf{型と契約が双方向に検証される。}

\paragraph{6.4.4
「メタ進化」のためのリフレクション}\label{ux30e1ux30bfux9032ux5316ux306eux305fux3081ux306eux30eaux30d5ux30ecux30afux30b7ux30e7ux30f3}

PhiLang は自身のフェーズを進化計算の対象にできる。本実験の
\texttt{.aice} (YAML) に書かれた 11 変数は、PhiLang で書けば
\texttt{evolvable} 修飾子付きパラメータになる:

\begin{Shaded}
\begin{Highlighting}[]
\NormalTok{phase Pi\_Phase1\_Robustness}
\NormalTok{  evolvable concurrency  in [c1, c2, c4, c8, c16, c32]}
\NormalTok{  evolvable seed\_count   in [s4, s6, s8, s12, s16, s24]}
\NormalTok{  evolvable generations  in [g10, g15, g20, g30, g40, g60, g80]}
\NormalTok{  ...}
\NormalTok{  optimize for \{}
\NormalTok{     maximize completed\_generations}
\NormalTok{     minimize cost\_usd}
\NormalTok{  \}}
\end{Highlighting}
\end{Shaded}

→ \textbf{PhiLang は MAP-Elites の対象であると同時に、MAP-Elites
を記述する言語}でもある。これは PiLang
が自身を進化させないのに対し、PhiLang 固有の自己参照性である。

\subsection{7.
限界と次のステップ}\label{ux9650ux754cux3068ux6b21ux306eux30b9ux30c6ux30c3ux30d7}

\subsubsection{7.1 reviewer が見落とした特性 (F3 と provider
不整合)}\label{reviewer-ux304cux898bux843dux3068ux3057ux305fux7279ux6027-f3-ux3068-provider-ux4e0dux6574ux5408}

§4.1 の通り、\texttt{checkpoint\_per\_gen=off} が top 5 で
5/5、\texttt{provider=mock} が top 5 で 3/5 を占めた。これは reviewer
の設計バグであり PhiLang の知見ではない。次イテレーション (v2)
で以下を直す:

\begin{enumerate}
\def\labelenumi{\arabic{enumi}.}
\tightlist
\item
  \textbf{F3 を hard floor 化}: \texttt{checkpoint\_per\_gen=off}
  の個体は最終スコアを × 0
\item
  \textbf{mock provider を最終スコア化前に除外}
\item
  \textbf{\texttt{provider\ ×\ model} coherence rule 追加}:
  \texttt{openai\ ⇒\ model\ ∈\ \{gpt-4o,\ gpt-4o-mini\}} 等
\end{enumerate}

\subsubsection{7.2 PhiLang
のコンパイラが未実装}\label{philang-ux306eux30b3ux30f3ux30d1ux30a4ux30e9ux304cux672aux5b9fux88c5}

§6 の PhiLang は\textbf{設計}段階に留まっている。実装するには:

\begin{enumerate}
\def\labelenumi{\arabic{enumi}.}
\tightlist
\item
  PhiLang フロントエンド (\texttt{phi\_parser.py})
  を書く。\texttt{phase\ /\ every\ /\ after\ /\ guarantees\ /\ evolvable}
  の 5 ブロックを AST に。
\item
  \texttt{phi\_codegen.py} で env 設定 + CLI フラグ + AIPL
  コードに脱糖。
\item
  \texttt{guarantees} 節を \texttt{aipl\_typeck.py}
  の追加チェックとして実装。
\end{enumerate}

工数概算: 1〜2 日。

\subsubsection{7.3 YAML lowerer
の欠落}\label{yaml-lowerer-ux306eux6b20ux843d}

\texttt{aice\_parser.lower()} が curly-brace \texttt{.aice}
用なので、YAML 形式の本実験 \texttt{.aice} は手作業で \texttt{.ga.json}
にマップした。再現性のためには \texttt{aice\_parser\_yaml.py}
を新規に書くか、既存 lowerer に dialect 検出を追加する必要がある。

\subsubsection{7.4 進化軸が現実の API
制約を知らない}\label{ux9032ux5316ux8ef8ux304cux73feux5b9fux306e-api-ux5236ux7d04ux3092ux77e5ux3089ux306aux3044}

\texttt{provider\ ×\ model} の組合せ (\texttt{openai\ に\ claude-haiku})
のような実在しない組合せを reviewer は弾けなかった。これは reviewer
の知識の限界。次は coherence rule を充実させるか、reviewer の persona
に「実在する provider×model 対の白リスト」を埋め込む。

\subsection{8. まとめ}\label{ux307eux3068ux3081}

{\def\LTcaptype{none} % do not increment counter
\begin{longtable}[]{@{}
  >{\raggedright\arraybackslash}p{(\linewidth - 2\tabcolsep) * \real{0.5000}}
  >{\raggedright\arraybackslash}p{(\linewidth - 2\tabcolsep) * \real{0.5000}}@{}}
\toprule\noalign{}
\begin{minipage}[b]{\linewidth}\raggedright
項目
\end{minipage} & \begin{minipage}[b]{\linewidth}\raggedright
結果
\end{minipage} \\
\midrule\noalign{}
\endhead
\bottomrule\noalign{}
\endlastfoot
案 1 (YAML) \texttt{.aice} 設計 & 11 軸 / 5 reviewer / 5 task / 4 plan
branch を 230 行で記述 \\
案 3 (Math) \texttt{.aice} 設計 (参考) & 同じ問題を decision variables +
objective + constraints で記述 (270 行) \\
\texttt{.ga.json} 生成 & 手作業マッピング (codegen 互換) \\
\texttt{.aipl} 生成 & 21 KB / 641 行、parse OK、smoke +
本番ともに完走 \\
AIPL 本番ラン & 33 分 / \$0.269 / 950 calls / 38 個体 /
silent\_death=0 \\
\textbf{進化計算が選んだ言語} & \textbf{PhiLang} ---
\texttt{Bounded\ burst\ ×\ Mandatory\ retry\ ×\ Generous\ drain\ ×\ Forced\ checkpoint}
を\textbf{言語の意味論}で強制 \\
副次成果 & reviewer 設計の弱点 (weight 不足 / 不在 coherence) を 3
点同定 \\
\end{longtable}
}

本実験は
\textbf{「メタ計算問題からプログラミング言語を進化計算で生み出す」}
という仮説を、PiLang (π 数値計算問題) と同じ
\texttt{.aice\ →\ .ga.json\ →\ .aipl}
パイプライン上で\textbf{端から端まで動作させた第 2 の実証}である。

PiLang
が「\textbf{特定の数値計算}を最短に書く言語」だったのに対し、PhiLang
は「\textbf{自分自身の進化計算実行}を最も安全に書く言語」である点で対照的だが、両者とも

\begin{itemize}
\tightlist
\item
  \textbf{MAP-Elites の個体群から共通する設計原理を抽出}し、
\item
  \textbf{その原理を意味論レベルで強制する 3〜4 行の最小構文}
  に結晶化する、
\end{itemize}

という同じ方法論を踏襲している。\texttt{aice-evolution-v2}
パイプラインは\textbf{問題ドメインに依存しない}普遍的な「言語設計支援ツール」として機能することが、今回
PiLang 以外のドメイン (メタ層) で初めて検証された。

\begin{center}\rule{0.5\linewidth}{0.5pt}\end{center}

\subsection{参考: 主要ファイル
(このディレクトリ内)}\label{ux53c2ux8003-ux4e3bux8981ux30d5ux30a1ux30a4ux30eb-ux3053ux306eux30c7ux30a3ux30ecux30afux30c8ux30eaux5185}

\begin{itemize}
\tightlist
\item
  問題定義: \texttt{Pi\_Phase1\_Robustness.aice} (YAML),
  \texttt{Pi\_Phase1\_Robustness.math.aice} (数理 参考)
\item
  スキーマ: \texttt{robustness\_paradigm.schema.json}
\item
  中間 IR: \texttt{Pi\_Phase1\_Robustness.ga.json}
\item
  AIPL プログラム: \texttt{Pi\_Phase1\_Robustness.aipl}
\item
  実行結果 lineage: \texttt{Pi\_Phase1\_Robustness.aipl\_lineage.json}
\item
  課金記録: \texttt{ai\_usage.json}
\item
  実行ログ: \texttt{Pi\_Phase1\_Robustness.log}
\end{itemize}

\subsection{参考: 親実験}\label{ux53c2ux8003-ux89aaux5b9fux9a13}

\begin{itemize}
\tightlist
\item
  \texttt{aice-pi-evolution/REPORT\_JA.md} --- 2026-05-15 の PiLang
  探索報告。本実験はその「同じ方法論を別ドメインに適用した第 2
  例」に相当する。
\end{itemize}

\end{document}
